% element: 10-s074:e03
% element: 10-s074:e04
\BeleskeID{10-s074}
Svaka DC granica važi uz testne uslove navedene u istom redu. Kod $V_{OH}$ i $V_{OL}$ obavezno poveži garantovani napon sa odgovarajućom izlaznom strujom i familijom 54 ili 74. Ulazni pragovi, ulazne struje i zaštitni napon nisu međusobno zamenljivi parametri.

Struja kratkog spoja nije dozvoljena trajna struja opterećenja. Napomena ograničava ispitivanje na najviše jedan kratko spojen izlaz i najviše jednu sekundu. Struja napajanja razlikuje se za visok i nizak izlaz, što utiče na disipaciju.

AC podaci važe na $25^\circ\mathrm C$, uz $V_{CC}=5\,\mathrm V$ i $C_L=15\,\mathrm{pF}$. Tipična i maksimalna kašnjenja razlikuju se; pri vremenskom proračunu mora se koristiti odgovarajuća kolona i smer promene izlaza.

Pri DC ispitivanjima $V_{IH}$ i $V_{IL}$ su garantovani nivoi za sve ulaze. Za zaštitnu diodu važi $V_{CC}=\mathrm{MIN}$ i $I_{IN}=-18\,\mathrm{mA}$. Visoki izlazni nivo meri se uz $V_{CC}=\mathrm{MIN}$ i $I_{OH}=\mathrm{MAX}$; ulazi su na $V_{IH}$ ili $V_{IL}$ prema tabeli istinitosti. Niski izlazni nivo ima iste minimalne uslove napajanja, sa ulazima na $V_{IL}$ ili $V_{IH}$ prema tabeli. Ulazne struje i struje kratkog spoja i napajanja ispituju se pri $V_{CC}=\mathrm{MAX}$.
